% GENERATED FILE. Edit canonical lessons, then run npm run generate:books.
% canonical-source-hash: fnv1a64:ed7f9c4a080d9e71
% canonical-lessons: TE-C126-batu, TE-C126-pandi, TE-C126-gadida, TE-C126-onte, TE-C126-mukham

\chapter{Duck, Pig, Donkey, Camel, Face}
\label{ch:te-duck-pig-donkey-camel-face}
\hlchaptermodality{126}

\begin{chapteropening}
\textbf{By the end of this chapter:} \emph{I can say a duck, a pig, a donkey, a camel and the face.}

Complete the last lesson of chapter 126: I can say a duck, a pig, a donkey, a camel and the face.
\end{chapteropening}

\section[\texorpdfstring{\te{బాతు} (bātu)}{బాతు (bātu)}]{\te{బాతు} (bātu) — a duck}

\label{lesson:TE-C126-batu}



\begin{quote}
Before the new one: say the Telugu for a butterfly, then the Telugu for a crow.
\end{quote}

\subsection*{You'll want to know: \te{బాతు}}
\textbf{\te{బాతు}} — \emph{bātu} — \textquotedblleft{}a duck\textquotedblright{}.

\begin{grammarlens}[title={What you've built}]
One more word for this chapter.
\end{grammarlens}

\subsection*{Your turn}
\begin{itemize}
\raggedright
  \item \emph{Say it:} \emph{bātu}
  \item \emph{Say it:} \emph{bātu}, once more
  \item \emph{Say it:} say \emph{bātu}
  \item \emph{Recall:} say the Telugu for a butterfly, then the Telugu for a crow, then say \emph{bātu} again
\end{itemize}

\begin{tcolorbox}[breakable,colback=teal!4,colframe=teal!35!black,title={Before you move on}]
What does \te{బాతు} mean? (\textquotedblleft{}A duck\textquotedblright{}.) Say it once more.
\end{tcolorbox}
\section[\texorpdfstring{\te{పంది} (pandi)}{పంది (pandi)}]{\te{పంది} (pandi) — a pig}

\label{lesson:TE-C126-pandi}



\begin{quote}
Before the new one: say the Telugu for a crow, then the Telugu for a duck.
\end{quote}

\subsection*{You'll want to know: \te{పంది}}
\textbf{\te{పంది}} — \emph{pandi} — \textquotedblleft{}a pig\textquotedblright{}.

\begin{grammarlens}[title={What you've built}]
One more word for this chapter.
\end{grammarlens}

\subsection*{Your turn}
\begin{itemize}
\raggedright
  \item \emph{Say it:} \emph{pandi}
  \item \emph{Say it:} \emph{pandi}, once more
  \item \emph{Say it:} say \emph{pandi}
  \item \emph{Recall:} say the Telugu for a crow, then the Telugu for a duck, then say \emph{pandi} again
\end{itemize}

\begin{tcolorbox}[breakable,colback=teal!4,colframe=teal!35!black,title={Before you move on}]
What does \te{పంది} mean? (\textquotedblleft{}A pig\textquotedblright{}.) Say it once more.
\end{tcolorbox}
\section[\texorpdfstring{\te{గాడిద} (gāḍida)}{గాడిద (gāḍida)}]{\te{గాడిద} (gāḍida) — a donkey}

\label{lesson:TE-C126-gadida}



\begin{quote}
Before the new one: say the Telugu for a duck, then the Telugu for a pig.
\end{quote}

\subsection*{You'll want to know: \te{గాడిద}}
\textbf{\te{గాడిద}} — \emph{gāḍida} — \textquotedblleft{}a donkey\textquotedblright{}.

\begin{grammarlens}[title={What you've built}]
One more word for this chapter.
\end{grammarlens}

\subsection*{Your turn}
\begin{itemize}
\raggedright
  \item \emph{Say it:} \emph{gāḍida}
  \item \emph{Say it:} \emph{gāḍida}, once more
  \item \emph{Say it:} say \emph{gāḍida}
  \item \emph{Recall:} say the Telugu for a duck, then the Telugu for a pig, then say \emph{gāḍida} again
\end{itemize}

\begin{tcolorbox}[breakable,colback=teal!4,colframe=teal!35!black,title={Before you move on}]
What does \te{గాడిద} mean? (\textquotedblleft{}A donkey\textquotedblright{}.) Say it once more.
\end{tcolorbox}
\section[\texorpdfstring{\te{ఒంటె} (oṇṭe)}{ఒంటె (oṇṭe)}]{\te{ఒంటె} (oṇṭe) — a camel}

\label{lesson:TE-C126-onte}



\begin{quote}
Before the new one: say the Telugu for a pig, then the Telugu for a donkey.
\end{quote}

\subsection*{You'll want to know: \te{ఒంటె}}
\textbf{\te{ఒంటె}} — \emph{oṇṭe} — \textquotedblleft{}a camel\textquotedblright{}.

\begin{grammarlens}[title={What you've built}]
One more word for this chapter.
\end{grammarlens}

\subsection*{Your turn}
\begin{itemize}
\raggedright
  \item \emph{Say it:} \emph{oṇṭe}
  \item \emph{Say it:} \emph{oṇṭe}, once more
  \item \emph{Say it:} say \emph{oṇṭe}
  \item \emph{Recall:} say the Telugu for a pig, then the Telugu for a donkey, then say \emph{oṇṭe} again
\end{itemize}

\begin{tcolorbox}[breakable,colback=teal!4,colframe=teal!35!black,title={Before you move on}]
What does \te{ఒంటె} mean? (\textquotedblleft{}A camel\textquotedblright{}.) Say it once more.
\end{tcolorbox}
\section[\texorpdfstring{\te{ముఖం} (mukhaṁ)}{ముఖం (mukhaṁ)}]{\te{ముఖం} (mukhaṁ) — the face}

\label{lesson:TE-C126-mukham}



\begin{quote}
Before the new one: say the Telugu for a donkey, then the Telugu for a camel.
\end{quote}

\subsection*{You'll want to know: \te{ముఖం}}
\textbf{\te{ముఖం}} — \emph{mukhaṁ} — \textquotedblleft{}the face\textquotedblright{}.

\begin{grammarlens}[title={What you've built}]
One more word for this chapter.
\end{grammarlens}

\subsection*{Your turn}
\begin{itemize}
\raggedright
  \item \emph{Say it:} \emph{mukhaṁ}
  \item \emph{Say it:} \emph{mukhaṁ}, once more
  \item \emph{Say it:} say \emph{mukhaṁ}
  \item \emph{Recall:} say the Telugu for a donkey, then the Telugu for a camel, then say \emph{mukhaṁ} again
\end{itemize}

\begin{tcolorbox}[breakable,colback=teal!4,colframe=teal!35!black,title={Before you move on}]
What does \te{ముఖం} mean? (\textquotedblleft{}The face\textquotedblright{}.) Say it once more.
\end{tcolorbox}
